% part: intuitionistic-logic
% chapter: soundness-completeness
% section: lindenbaum

\documentclass[../../../include/open-logic-section]{subfiles}

\begin{document}

\olfileid{int}{sc}{lin}

\olsection{د لينډنباوم لمه}

د شهودي منطق د بشپړتيا قضيه د دې په فرضولو ثابتېږي چې
$\Gamma \Proves/ !A$؛ بيا داسې مدل جوړوو چې
$\mSat{M}{\Gamma}$ او~$\mSat/{M}{!A}$ ولري۔

په کلاسيک منطق کښې د !!{derivability} اړيکه سازګارۍ ته راکمولے
شو، ځکه !!a{formula}~$!A$ له د !!{formula}s له يوه سټ څخه هله او
يوازې هله !!{derivable} وي چې دغه سټ د~$!A$ له نفي سره يوځاے
ناسازګار وي۔ په شهودي منطق کښې دا کار نۀ شو کولے۔ په شهودي منطق
کښې که $\lnot!A$ ناسازګار وي، يوازې $\Proves \lnot\lnot !A$
ترلاسه کوو۔ څرنګه چې $\lnot\lnot!A \lif !A$ په عمومي ډول په
شهودي منطق کښې نۀ منل کېږي، نو $\Proves !A$ نۀ شو اخيستلے۔

نو د مدل~$\mModel{M}$ د جوړولو پر مهال بايد د
!!{formula}~$!A$ د !!{derivability} نشتون په پام کښې وساتو؛ له دې امله د مدل
$\mModel{M}$ د جوړولو دپاره بشپړ سټ $\Gamma^* \supseteq \Gamma$
نۀ شو کارولے، ځکه په هر بشپړ سټ~$\Gamma^*$ کښې
$\Gamma^* \Proves !A \lor \lnot !A$ وي۔

د بشپړ سټ~$\Gamma^*$ پر ځاے د فارمولونو د اولني سټ مفهوم کاروو:

\begin{defn}\ollabel{defn:prime}
  د !!{formula}s سټ~$\Gamma$ هله او يوازې هله \emph{اولنی} دے چې
  \begin{enumerate}
  \item\ollabel{defn:prime1} $\Gamma$ سازګار وي، يعنې $\Gamma \Proves/ \lfalse$؛
  \item\ollabel{defn:prime2} که $\Gamma \Proves !A$ وي، نو $!A \in \Gamma$ وي؛ او
  \item\ollabel{defn:prime3} که $!A \lor !B \in \Gamma$ وي، نو $!A \in
    \Gamma$ يا $!B \in \Gamma$ وي۔
  \end{enumerate}
\end{defn}

\begin{lem}[د لينډنباوم لمه]
  \ollabel{lem:lindenbaum} که $\Gamma \Proves/ !A$ وي، نو داسې
  $\Gamma^* \supseteq \Gamma$ شته چې $\Gamma^*$ اولنی وي او
  $\Gamma^* \Proves/ !A$ وي۔
\end{lem}

\begin{proof}
  $!B_1 \lor !C_1$، $!B_2 \lor !C_2$، \dots، دې د
  $!B \lor !C$ په بڼه د ټولو !!{formula}s يوه شمېرنه وي۔ د
  !!{formula}s د سټونو يوه زياتېدونکې لړۍ~$\Gamma_n$ تعريفوو،
  چې هر $\Gamma_{n+1}$ د $\Gamma_n$ او يوه نوي !!{formula}
  له يوځاے کولو جوړېږي۔ $\Gamma^*$ به د ټولو~$\Gamma_n$ اتحاد
  وي۔ نوي !!{formula}s داسې ټاکو چې $\Gamma^*$ اولنی شي او لا
  هم $\Gamma^* \Proves/ !A$ پاتې شي۔ يعنې په هر ګام کښې بايد
  لومړنی داسې فصل~$!B_i \lor !C_i$ ومومو چې:
  \begin{enumerate}
  \item\ollabel{gamma-1} $\Gamma_n \Proves !B_i \lor !C_i$
  \item\ollabel{gamma-2} $!B_i \notin \Gamma_n$ او $!C_i \notin \Gamma_n$
  \end{enumerate}
  يا $!B_i$ په $\Gamma_n$ کښې زياتوو، که $\Gamma_n \cup \{!B_i\}
  \Proves/ !A$ وي؛ او که داسې نۀ وي، $!C_i$ زياتوو۔ بايد وښيو
  چې دا کار سم دے۔ اوس $i(n)$ هغه تر ټولو کوچنی $i$ ټاکو چې
  \olref{gamma-1} او~\olref{gamma-2} پوره کوي۔
  
  $\Gamma_0 = \Gamma$ او
  \[
  \Gamma_{n+1} = \begin{cases}
    \Gamma_n \cup \{!B_{i(n)}\} &
    \text{که $\Gamma_n \cup \{!B_{i(n)}\} \Proves/ !A$ وي} \\
    \Gamma_n \cup \{!C_{i(n)}\} & \text{که داسې نۀ وي}
    \end{cases}
  \]
  تعريف کړئ۔ که $i(n)$ نۀ وي ټاکلے، يعنې هر کله چې
  $\Gamma_n \Proves !B \lor !C$ وي، نو يا $!B \in \Gamma_n$
  يا $!C \in \Gamma_n$ وي، $\Gamma_{n+1}
  = \Gamma_n$ ږدو۔ اوس $\Gamma^* = \bigcup_{n=0}^\infty \Gamma_n$
  ږدو۔

  لومړے ښيو چې د هر $n$ دپاره $\Gamma_n \Proves/ !A$ وي۔ پر~$n$
  استقرا کاروو۔ د $n = 0$ لپاره دعوه د لمې له فرضه، يعنې
  $\Gamma \Proves/ !A$، راځي۔ که $n>0$ وي، بايد وښيو چې له
  $\Gamma_n \Proves/ !A$ څخه $\Gamma_{n+1} \Proves/ !A$
  راځي۔ که $i(n)$ نۀ وي ټاکلے، $\Gamma_{n+1} = \Gamma_n$ او
  د ثابتولو څۀ نشته۔ نو فرض کړئ $i(n)$ ټاکلے وي۔ د اسانتيا دپاره
  $i = i(n)$ ږدو۔
  
  د دعوې مقابل لوری ثابتوو۔ فرض کړئ $\Gamma_{n+1}
  \Proves !A$۔ د جوړونې له مخې $\Gamma_{n+1} = \Gamma_n \cup
  \{!B_i\}$ وي که $\Gamma_n \cup \{!B_i\} \Proves/ !A$ وي؛
  او که داسې نۀ وي، $\Gamma_{n+1} = \Gamma_n \cup \{!C_i\}$
  وي۔ لومړے حالت څرګند ناشونی دے، ځکه په هغه کښې به
  $\Gamma_{n+1} \Proves/ !A$ وي۔ نو $\Gamma_n \cup
  \{!B_i\} \Proves !A$ او $\Gamma_{n+1} = \Gamma_n \cup
  \{!C_i\}$ دي۔ د $i(n)$ له تعريفه $\Gamma _n \Proves !B_i \lor
  !C_i$ لرو۔ $\Gamma_n \cup \{!B_i\} \Proves !A$ هم لرو، او
  $\Gamma_{n+1} = \Gamma_n \cup \{!C_i\} \Proves !A$ هم۔ نو
  $\Gamma_n \Proves !A$ راځي، لکه چې غوښتل مو وښيو۔

  که $\Gamma^* \Proves !A$ وي، نو داسې متناهي فرعي سټ
  $\Gamma' \subseteq \Gamma^*$ شته چې $\Gamma' \Proves !A$
  وي۔ هر $!D \in \Gamma'$ بايد د کوم~$i$ دپاره په $\Gamma_i$
  کښې وي۔ $n$ دې د دغو شاخصونو تر ټولو لوے وي۔ څرنګه چې د
  $i \le n$ لپاره $\Gamma_i \subseteq \Gamma_n$ وي، نو
  $\Gamma' \subseteq \Gamma_n$۔ خو بيا $\Gamma_n \Proves !A$
  کېږي؛ دا زموږ له پورته ثبوت شوي $\Gamma_n \Proves/ !A$ سره
  په ټکر کښې دے۔

  په پاے کښې ښيو چې $\Gamma^*$ اولنی دے، يعنې د
  \olref{defn:prime} درې شرطونه \olref{defn:prime1}،
  \olref{defn:prime2} او~\olref{defn:prime3} پوره کوي۔

  لومړے، $\Gamma^* \Proves/ !A$ دے، نو $\Gamma^*$ سازګار
  دے؛ ځکه \olref{defn:prime1} پوره کېږي۔

  اوس ښيو چې که $\Gamma^* \Proves !B \lor !C$ وي، نو يا
  $!B \in \Gamma^*$ يا $!C \in \Gamma^*$ وي۔ په دې سره
  \olref{defn:prime3} ثابتېږي، ځکه که $!B \lor !C \in \Gamma^*$
  وي، نو $\Gamma^* \Proves !B \lor !C$ هم وي۔ نو فرض کړئ
  $\Gamma^* \Proves !B \lor !C$ وي، خو $!B \notin \Gamma^*$
  او $!C \notin \Gamma^*$ وي۔ له $\Gamma^* \Proves !B \lor !C$
  څخه د کوم~$n$ دپاره $\Gamma_n \Proves !B \lor !C$ راځي۔
  $!B \lor !C$ د ټولو فصلونو په شمېرنه کښې، ووايئ د
  $!B_j \lor !C_j$ په توګه، راځي۔ دغه $!B_j \lor !C_j$ د
  $i(n)$ د تعريف شرطونه پوره کوي: $\Gamma_n \Proves !B_j
  \lor !C_j$ او $!B_j \notin \Gamma_n$ او $!C_j \notin \Gamma_n$
  دي۔ د دغه شاخص څخه ټيټ شاخصونه متناهي دي؛ هر ټيټ فصل
  $!B_i \lor !C_i$ چې په کوم پړاو کښې غوره شي، د يوه غړي
  $!B_i$ يا $!C_i$ په زياتېدو سره بيا شرط نۀ پوره کوي۔
  ځکه چې دغه فصل د فرض له مخې لا شرط پوره کوي، په يو پړاو~$m$
  کښې به $j = i(m)$ شي۔ نو يا $!B \in \Gamma_{m+1}$ يا
  $!C \in \Gamma_{m+1}$ کېږي، چې د $!B \notin \Gamma^*$ او
  $!C \notin \Gamma^*$ له فرض سره ټکر لري۔ د ژباړې د نسخې
  څرګندونه (OLINT-009): منجمده سرچينه دلته وايي چې په هر ګام
  کښې د شرط پوره کوونکو فصلونو ټول شمېر کمېږي؛ د نوي غړي په
  زياتېدو سره نور فصلونه هم ثابتېدے شي۔ له ټاکلي شاخص څخه ټيټ
  متناهي شاخصونه يو ځل هوارېږي، او همدا د دې دليل کافي بڼه ده۔

  اوس فرض کړئ $\Gamma^* \Proves !B$۔ نو $\Gamma^* \Proves !B \lor
  !B$ وي۔ خو تازه مو ثابت کړل چې که $\Gamma^* \Proves !B \lor !B$
  وي، نو $!B \in \Gamma^*$ وي۔ نو $\Gamma^*$ د
  \olref{defn:prime} شرط~\olref{defn:prime2} پوره کوي۔
\end{proof}

\begin{prob}
وښيئ چې که $\Gamma \Proves/ \bot$ وي، نو $\Gamma$ په کلاسيک
منطق کښې سازګار دے؛ يعنې داسې ټاکنه شته چې د~$\Gamma$ ټول
فارمولونه رښتيا کوي۔
\end{prob}

\end{document}
